\editionPage[\pdfbookmark[0]{Démonstrations}{bm-preuves-debut-parent}]{preuves-debut}{Comprendre les preuves}{Pourquoi le produit de deux négatifs est-il positif ?}{Les arguments ci-dessous éclairent des propriétés utilisées dans le livre.}
\begin{blocEdition}{preuves-debut}

\editionEtape{Partir de l’opposé}Soient $a$ et $b$ deux nombres positifs. Comme $b+(-b)=0$, multiplier par $a$ donne $ab+a(-b)=0$ par distributivité. Ainsi $a(-b)$ est l’opposé de $ab$ : $a(-b)=-ab$.
\editionEtape{Changer aussi le premier signe}De même, $a+(-a)=0$. En multipliant par $-b$, on obtient $a(-b)+(-a)(-b)=0$. On sait que $a(-b)=-ab$ ; donc $(-a)(-b)$ est l’opposé de $-ab$, c’est-à-dire $ab$.
\editionEtape{Ce que prouve le raisonnement}La règle des signes est imposée par la distributivité et par la notion d’opposé : ce n’est pas une convention isolée à mémoriser. Le cas d’un facteur nul se traite à part : le produit vaut 0.
\editionRetenir{Deux facteurs négatifs donnent un produit positif, parce que ce produit doit être l’opposé d’un produit négatif pour que la distributivité reste vraie.}
\end{blocEdition}
\editionPage{preuves-2}{Comprendre les preuves}{Pourquoi diviser par une fraction revient-il à multiplier par son inverse ?}{Les arguments ci-dessous éclairent des propriétés utilisées dans le livre.}
\begin{blocEdition}{preuves-2}

\editionEtape{Utiliser la définition du quotient}Pour $v\ne0$, le quotient $u:v$ est l’unique nombre $q$ tel que $q\times v=u$.
\editionEtape{Tester une expression}L’inverse de $v$ vérifie $v\times\frac1v=1$. Le nombre $u\times\frac1v$, multiplié par $v$, donne $u\times\left(\frac1v\times v\right)=u$. C’est donc le quotient $u:v$.
\editionEtape{Appliquer aux fractions}Si $v=\frac cd$ avec $c$ et $d$ non nuls, son inverse est $\frac dc$, car $\frac cd\times\frac dc=1$. Ainsi $\frac ab:\frac cd=\frac ab\times\frac dc$, tous les dénominateurs étant non nuls.
\editionEtape{Un exemple contrôlé}$\frac34:\frac25=\frac34\times\frac52=\frac{15}{8}$. En sens inverse : $\frac{15}{8}\times\frac25=\frac34$. On inverse uniquement le diviseur.
\editionRetenir{Pour contrôler un quotient, remultiplier le résultat par le diviseur : on doit retrouver le dividende.}
\end{blocEdition}
\editionPage{preuves-3}{Comprendre les preuves}{Le théorème de Pythagore par une comparaison d’aires}{Les arguments ci-dessous éclairent des propriétés utilisées dans le livre.}
\begin{blocEdition}{preuves-3}

\editionEtape{Construire la figure}On place quatre triangles rectangles identiques, de côtés de l’angle droit $a$ et $b$ et d’hypoténuse $c$, dans un carré de côté $a+b$. Le quadrilatère central a quatre côtés de longueur $c$. Ses angles sont droits, car les deux angles aigus d’un triangle rectangle ont pour somme $90^\circ$ : c’est un carré.
\label{fig-preuves-3-preuve-pythagore}\figureGeo{preuve-pythagore}

\editionEtape{Compter deux fois la même aire}Le grand carré a pour aire $(a+b)\times(a+b)$. La somme des aires des quatre triangles et du carré central vaut $4\times\frac{ab}2+c^2=2ab+c^2$.
\editionEtape{Comparer}En découpant autrement le carré de côté $a+b$, en deux carrés de côtés $a$ et $b$ et deux rectangles de côtés $a$ et $b$, son aire vaut aussi $a^2+b^2+2ab$. Donc $a^2+b^2+2ab=2ab+c^2$, et en retirant $2ab$ : $\res{a^2+b^2=c^2}$.
\editionRetenir{Le théorème compare les aires des carrés construits sur les côtés. La preuve utilise l’hypothèse du triangle rectangle : elle est essentielle.}
\end{blocEdition}
\editionPage{preuves-4}{Comprendre les preuves}{Les trois théorèmes de la droite des milieux}{Les arguments ci-dessous éclairent des propriétés utilisées dans le livre.}
\begin{blocEdition}{preuves-4}

\editionEtape{Un premier parallélogramme}Dans le triangle ABC, I est le milieu de $[AB]$ et J celui de $[AC]$. Le demi-tour de centre J transforme I en K : J est le milieu de $[IK]$. Les diagonales $[AC]$ et $[IK]$ du quadrilatère AICK ont le même milieu J : AICK est un parallélogramme. Donc $(CK)\parallel(AB)$ et $CK=AI$.
\label{fig-preuves-4-preuve-milieux}\figureGeo{preuve-milieux}

\editionEtape{Un second parallélogramme}Comme I est le milieu de $[AB]$, $AI=IB$, donc $CK=IB$ ; de plus $(CK)\parallel(IB)$. Le quadrilatère non croisé IBCK a deux côtés opposés parallèles et de même longueur : c’est un parallélogramme. Ainsi $(IK)\parallel(BC)$ et $IK=BC$. Comme J est le milieu de $[IK]$ : $\res{(IJ)\parallel(BC)}$ et $\res{IJ=\frac{BC}2}$. Ce sont les théorèmes 1 et 2.
\editionEtape{Le troisième théorème}La parallèle $(d)$ à $(BC)$ passant par I coupe $(AC)$ en un point M. D’après le théorème 1, la droite $(IJ)$ est elle aussi parallèle à $(BC)$ et passe par I. Par un point, il ne passe qu’une seule parallèle à une droite donnée : $(IJ)=(d)$. Les droites $(d)$ et $(AC)$ n’ayant qu’un point commun, $M=J$ : $(d)$ coupe $[AC]$ en son milieu.
\editionRetenir{Les théorèmes de la droite des milieux reposent sur les propriétés du demi-tour et du parallélogramme vues en Cinquième. Le troisième utilise l’unicité de la parallèle.}
\end{blocEdition}
